\chapter{Introduction}
\label{chap:introduction}

Engineering a collective system means deciding three things:
what the collective computes,
where each piece of that computation runs,
and what to do with that placement once the infrastructure has changed.
%
The three decisions are taken by unrelated means, at different moments:
the computation in a programming model,
the placement in a deployment descriptor written before the system starts,
and its revision by an operator working from outside the application.

What a \ac{CAS} computes is a property of the collective:
the system is built from many autonomous entities that interact locally,
and its useful behaviour belongs to the collective they form and not to any single one of them~\cite{DBLP:journals/sttt/NicolaJW20}.
%
Its membership is not fixed:
entities join, leave, and fail while the system runs,
and an ensemble is identified by the properties its members hold at the time~\cite{DBLP:journals/taas/NicolaLPT14}.
%
A specification therefore states what the collective achieves, and what any single entity does follows from it:
a robot team holds a formation~\cite{DBLP:journals/swarm/BrambillaFBD13},
the phones carried through a public event estimate how crowded it is~\cite{DBLP:journals/computer/BealPV15}.
%
Correctness is argued the same way, over the behaviour of the collective as membership and environment change;
self-stabilisation is one established criterion~\cite{DBLP:journals/tomacs/ViroliABDP18}.

Such a system can take advantage of being deployed across the edge-cloud continuum,
where computation can sit close to the devices that produce its data,
or on hosts with the capacity to run it.
%
The same infrastructure works against it:
hosts differ by orders of magnitude in capability,
join and leave without notice,
and lose connectivity to one another.
%
The program does not say which host runs which piece of it, so someone has to decide.

That decision is normally taken once and left untouched.
%
This thesis makes it a variable of the running system.
%
A macro-program, the program that states what the collective achieves,
can be partitioned into components whose collective behaviour does not depend on where each component runs;
under conditions this thesis states and proves,
the assignment of components to hosts stops being a correctness concern.
%
It becomes a choice that can be made and remade while the system runs,
on non-functional grounds alone: battery, latency, load, carbon.

The model alone does not build a system.
%
The components it partitions have to come from a program,
and today such a program is written in one of three paradigms: aggregate, choreographic, or multitier.
%
Each yields deployable units in terms that do not translate into the others.
%
Choosing a deployment and revising it as the infrastructure drifts then takes a middleware that can relocate a running component,
and a policy that says which deployment to move to.
%
This thesis develops both sides: the language support, and the platform with the policies that drive it.

\section{Context and Motivation}
\label{sec:context-and-motivation}

Latency budgets, data locality, bandwidth, and intermittent connectivity pushed computation out of the cloud~\cite{DBLP:journals/iotj/ShiCZLX16,DBLP:journals/computer/Satyanarayanan17}.
%
What remains is a continuum of hosts running from the microcontroller embedded in the physical world to the server rack in the data centre (\cref{sec:the-continuum,subsec:continuum-characteristics}).
%
Its topology is unstable: devices move, batteries drain, links fail, and the network partitions (\cref{subsec:network-volatility}).
%
A placement that was good when the system started need not stay good for an hour.

This thesis targets two families of collective system on that infrastructure:
\ac{IoT} ecosystems~\cite{DBLP:journals/cn/AtzoriIM10} and robot swarms~\cite{DBLP:journals/swarm/BrambillaFBD13} (\cref{sec:target-environments}).
%
Both are built from many devices, each unreliable and each replaceable.
%
A correctness argument assembled device by device therefore has to be reassembled whenever a device joins or leaves~\cite{DBLP:journals/sttt/NicolaJW20}.

Macroprogramming moves the unit of description to the collective:
one program states what the collective achieves, and each device's behaviour is derived from it~\cite{DBLP:journals/csur/Casadei23} (\cref{sec:macro-programming}).
%
Three paradigms realise that idea (\cref{sec:foundational-paradigms}).
%
Aggregate computing computes over a field spanning an ensemble;
choreographic programming states a protocol among roles and projects it onto each participant;
multitier programming annotates the scopes of a single program with the tier they belong to.
%
All three derive logical units of computation from a global specification.

None of them decides where those units run.
%
Aggregate computing assumes every device evaluates a complete instance of the program, never a fragment;
endpoint projection fixes one process per role;
multitier placement is resolved at compile time.
%
The mapping onto physical hosts is therefore either implicit in the execution model or settled before the program is released,
and it is that mapping a volatile infrastructure keeps invalidating.

The architectures used to structure distributed software do not let the application set that mapping while it runs either (\cref{sec:architectural-paradigms}).
%
A microservice's deployment unit is a container under an orchestrator,
and re-placing it is an operations action that the application can neither express nor trigger.
%
\ac{FaaS} narrows the unit but makes it stateless,
so application state moves to backing services, usually cloud-hosted.
%
Actor runtimes keep state inside the actor and can place actors automatically,
but the placement policy belongs to the runtime and is tuned for data-centre clusters (\cref{subsec:paradigm-limitations}).

Pulverisation closes part of that gap.
%
Casadei et al.\ split each logical device into sensing, actuation, behaviour, state, and communication components,
each deployable wherever the continuum offers the resources it needs~\cite{DBLP:journals/fi/CasadeiPPVW20}.
%
The partitioning is fixed in advance, though:
five deployable units per logical device whatever the application computes,
and a logical device is relocated as a whole.
%
So is the deployment: the model says which deployments are admissible,
not who chooses among them or on what grounds (\cref{chap:dynamic-reconfiguration}).
%
PulvReAKt, an earlier step of this work, moved that assignment into a declarative specification and made it revisable at runtime~\cite{DBLP:journals/fgcs/FarabegoliPCV24}.
% the five roles stayed as they were, and it still took a person to say which deployment to move to.

\section{Research Questions}
\label{sec:research-questions}

Three questions follow from that gap.
%
The model comes first because it licenses the other two:
until a deployable component is defined, and the consequences of moving one are known,
the mapping of components to hosts cannot be posed as a runtime decision.

\rqhead{rq:model}{Deployment model}
A macro-program is written against the collective,
and its components end up on devices that differ in what they can run.
%
The five-role partitioning fixes how many deployable units there are before the application is known,
and the paradigms that produce those units leave their placement implicit.
%
\emph{Under what conditions can a collective macro-program be partitioned into independently deployable components so that its observable behaviour does not depend on how those components are mapped onto physical devices?}
%
Answering it requires stating the conditions and proving the invariance under them.
%
Both are given in \cref{chap:pulverization}, for macro-programs partitioned along their own data flow (\cref{subsec:macroprogram-dag}),
and simulation measures how far apart two deployments run before they agree (\cref{subsec:pulverization-validation}).
%
A negative answer would be a well-formed macro-program with two correct deployments whose outputs settle on different values.

\rqhead{rq:languages}{Language and coordination support}
The three paradigms each produce deployable components, in terms that do not translate into one another:
a placed value, a projected role, and a field over an ensemble are not the same kind of object.
%
A system built from more than one of them therefore cannot be pulverised as a whole,
and running separate frameworks side by side gives up the guarantees each provides on its own.
%
\emph{What linguistic support lets a single program combine aggregate, choreographic, and multitier coordination -- and yield the deployable components the model requires -- while statically rejecting the coordination errors their ad-hoc combination admits?}
%
Two requirements pull against each other:
the discipline has to be sound, and it has to keep the programs the three paradigms support writable.
%
A soundness proof settles the first (\cref{subsec:camil-formalization}),
and reimplementing use cases drawn from the literature of all three settles the second (\cref{subsec:camil-expressivity}).
%
A negative answer would be a discipline that accepts a program violating one paradigm's own rules,
or one so restrictive that those use cases can no longer be written.

\subrqhead{rq:generation}{Generation}
Whatever the discipline, a program still has to be written,
in a \ac{DSL} whose operators are rare in the corpora general-purpose models are pretrained on.
%
\emph{Can such a program be obtained from a natural-language specification rather than written by hand, and what must a model be told about the language for that to work?}
%
The evidence is generated programs scored with \emph{pass@k} against test cases co-designed with the tasks and checked in simulation (\cref{sec:llm-pipeline}),
over sixteen models and three levels of supplied knowledge (\cref{sec:llm-evaluation}).
%
A negative answer would be models that stay at zero however much they are told,
or an outcome decided by model size and not by what the model is given.

\rqhead{rq:deployment}{Runtime deployment decision}
Deployment independence makes the choice among admissible deployments free, and leaves it unmade.
%
Whatever is chosen, the infrastructure will invalidate it,
so the choice has to be revisable while the system runs (\cref{sec:the-deployment-problem}).
%
\emph{Once the choice among admissible deployments is functionally free, how can that choice be made and revised while the system runs, on non-functional grounds alone -- by a policy stated in advance or by one learned from experience -- and what does the improvement in the targeted quantity cost in the others?}
%
The second half of the question is what makes it answerable:
showing that a policy improves what it targets settles nothing on its own,
so the simulations measure the cost elsewhere as well (\cref{chap:deployment-and-reconfiguration,chap:heterogeneous-gnn}).
%
A negative answer would be a policy whose gain in one quantity is paid for by an unbounded loss in another,
or one that reaches its target only by changing what the application computes.

\section{Contributions}
\label{sec:contributions}

Each contribution below names the chapter that develops it, the publication it comes from, and the question it answers;
\cref{tab:contributions} collects the same mapping.

\paragraph{A deployment model for macro-programs.}
A macro-program is partitioned along its own data flow into a \ac{DAG} of components:
one deployable unit per component of that flow,
each classified by whether its computation involves neighbouring devices,
and each reached through a forwarding chain from the device that owns it (\cref{sec:pulverization-abstractions}).
%
Under two hypotheses -- every component computes a self-stabilising function,
and the sensor state is uniform along every forwarding chain --
the behaviour observed at the macro-program's output ports stabilises to a value that does not depend on the deployment (\cref{thm:deployment-independence}).
%
Simulations in Alchemist over networks of 50, 100, and 200 mobile nodes compare a monolithic deployment against an offloaded one,
and the two converge to the same result once the nodes stop moving (\cref{subsec:pulverization-validation}).
%
The partitioning and the theorem are the thesis' central technical result, published as~\cite{DBLP:conf/acsos/FarabegoliVC24}, and they answer \cref{rq:model}.
%
An earlier step of this work, the PulvReAKt middleware, made the component-to-host assignment revisable at runtime through a declarative specification,
while leaving the original five-role partitioning unchanged~\cite{DBLP:journals/fgcs/FarabegoliPCV24}.
%
The two partitionings cut along different axes and compose (\cref{sec:pulverization-discussion}).

\paragraph{One type discipline for three coordination paradigms.}
\camil turns each paradigm into a \emph{capability}: a value that code must hold in order to use that paradigm's operations.
%
Placement types are shared by all of them,
so a value produced under one paradigm reaches another through its placement, with no ad-hoc bridge between them (\cref{sec:capabilities-unification}).
%
The discipline is realised with the contextual abstractions, capture checking, and separation checking of Scala~3, and it is proved sound (\cref{thm:camil-soundness}).
%
The evaluation reimplements 46 use cases drawn from the multitier, choreographic, and aggregate literature, with none left out,
and three applications combine paradigms within a single program (\cref{sec:camil-evaluation}).
%
The design is under review as~\cite{Farabegoli:2026:CaMiL}, and with the next contribution it answers \cref{rq:languages}.

\paragraph{Differentiated communication primitives for multiparty coordination.}
Choreographic and multitier languages offer point-to-point exchange and leave every other pattern to be encoded on top of it,
usually by broadcasting and filtering.
%
ScalaTropy gives each pattern an operation of its own:
isotropic, anisotropic, and co-anisotropic communication alongside the point-to-point case.
%
The \ac{DSL} is monadic, and its types record the architecture a choreography is written against (\cref{sec:scalatropy}).
%
On a matrix-vector product, the selective primitive keeps the bytes the master transmits nearly flat as workers are added,
while the same scatter encoded as broadcast-and-filter grows linearly (\cref{subsec:scalatropy-evaluation}).
%
The design is accepted for publication as~\cite{Farabegoli:2026:ScalaTropy}.

\paragraph{A body of knowledge for generating aggregate programs.}
An \ac{LLM} handles the surface syntax of a Scala \ac{DSL} fluently, but misses the meaning of the operators that make it a field language.
%
The \emph{body of knowledge} supplies that meaning:
a short operational manual, placed in the model's context at generation time,
introducing each operator through a minimal program and the values it produces over successive rounds (\cref{sec:body-of-knowledge}).
%
Its entries are added only where an observed failure demanded one, and never where an entry would contain a test case's solution (\cref{subsec:bok-codesign}).
%
Across sixteen models, the outcome is decided by what the model is told and not by model scale:
the largest model in the study leads at no knowledge level (\cref{sec:llm-evaluation}).
%
The contribution is published as~\cite{DBLP:journals/tiot/AguzziFV25} and answers \cref{rq:generation}.

\paragraph{Three hand-written reconfiguration policies.}
A device offloads one expensive component as soon as its own battery falls below a threshold, and moves it again when the host becomes unreachable~\cite{DBLP:conf/acsos/FarabegoliVC24}.
%
Infrastructural devices compete for the components application devices are better off not running, and each resizes the region it leads through a collectively built field~\cite{DBLP:journals/iot/FarabegoliPCV24}.
%
A planner searches the admissible deployments against a stated objective over energy and carbon, and repeats the search as the network drifts~\cite{DBLP:conf/coordination/BrogiCFFV25}.
%
Each policy reduces the quantity it targets at a bounded cost in another (extra messages, a stabilisation transient, or added latency),
and each requires a person to state in advance what a good deployment is (\cref{chap:deployment-and-reconfiguration}).
%
The three answer the first half of \cref{rq:deployment}.

\paragraph{A learned policy for collective offloading.}
The continuum is represented as a heterogeneous graph, with application devices, edge servers, and cloud servers as distinct node types,
and a \ac{GNN} over that graph serves as the Q-function of a deep Q-learning agent (\cref{sec:hetero-graph-state}).
%
An aggregate program running on the same devices whose deployment is being decided supplies a collective summary of local density that the graph alone does not carry (\cref{subsec:collective-information}).
%
Under a scalarised reward the learned policies move continuously between the extremes the weights ask for,
and with the collective term they differentiate by position, offloading part of a device's components at the fringe of a crowd, which no binary rule can do (\cref{sec:idhgql-evaluation}).
%
This removes the requirement that a person state the rule, and so answers the second half of \cref{rq:deployment}.
%
It is published as~\cite{Farabegoli:2026:HeteroGNN}.

\paragraph{A physical demonstrator.}
Project Emerge runs a single aggregate program on nine mobile robots that form spatial patterns and recover them after one of them is displaced by hand (\cref{chap:demonstrator}).
%
The robots carry a microcontroller and no aggregate runtime, and a server evaluates the program on their behalf,
which in the terms of the model is a deployment with the robots as owners and the server as their shared surrogate (\cref{sec:emerge-as-pulverisation}).
%
This is scoped realisability evidence for \cref{rq:model}:
a deployment the model admits, running on hardware and not in simulation.
%
It exercises none of the reconfiguration policies and reports no measurements, so it is not end-to-end validation of the thesis (\cref{sec:emerge-discussion}).
%
The contribution is published as~\cite{Farabegoli:2027:ProjectEmerge}, extending~\cite{DBLP:conf/coordination/AguzziBBCCDFPV25}.

\begin{table}[t]
    \centering
    \small
    \caption[Contributions, chapters, publications, and research questions]{Each contribution, the chapter that develops it, the publication it comes from, and the research question it answers.}
    \label{tab:contributions}
    \setlength{\tabcolsep}{4pt}
    \begin{tabularx}{\textwidth}{@{}>{\raggedright\arraybackslash}Xcll@{}}
        \toprule
        \textbf{Contribution} & \textbf{Ch.} & \textbf{Publication} & \textbf{Question} \\
        \midrule
        Declarative runtime reconfiguration of a pulverised deployment (PulvReAKt)
            & \labelcref{chap:pulverization} & FGCS 2024 & \cref{rq:model} \\
        Data-flow partitioning of a macro-program, with deployment independence
            & \labelcref{chap:pulverization} & ACSOS 2024 & \cref{rq:model} \\
        \camil: capabilities unifying three coordination paradigms
            & \labelcref{chap:language-and-coordination} & TOPLAS, under review & \cref{rq:languages} \\
        ScalaTropy: monadic communication primitives
            & \labelcref{chap:language-and-coordination} & COORDINATION 2026 & \cref{rq:languages} \\
        A body of knowledge for generating aggregate programs
            & \labelcref{chap:language-based-macroprogramming} & ACM TIoT 2025 & \cref{rq:generation} \\
        Battery-triggered offloading rule
            & \labelcref{chap:deployment-and-reconfiguration} & ACSOS 2024 & \cref{rq:deployment} \\
        Field-based reconfiguration over self-organising coordination regions
            & \labelcref{chap:deployment-and-reconfiguration} & Internet of Things 2024 & \cref{rq:deployment} \\
        Declarative planner for green deployments
            & \labelcref{chap:deployment-and-reconfiguration} & COORDINATION 2025 & \cref{rq:deployment} \\
        Learned collective offloading over a heterogeneous graph
            & \labelcref{chap:heterogeneous-gnn} & FGCS 2026 & \cref{rq:deployment} \\
        Project Emerge: one admissible deployment on physical robots
            & \labelcref{chap:demonstrator} & SCP 2027 & \cref{rq:model} \\
        \bottomrule
    \end{tabularx}
\end{table}

\section{Structure of the Thesis}
\label{sec:thesis-structure}

Part~I is background:
the continuum and the architectural paradigms used on it (\cref{chap:edge-cloud-continuum}),
the macroprogramming paradigms and the linguistic machinery built around them (\cref{chap:programming-models}),
and the placement, orchestration, and learning techniques a deployment decision can draw on (\cref{chap:dynamic-reconfiguration}).
%
Part~II states the deployment model and then the language support that feeds it,
from the model itself (\cref{chap:pulverization}) to the two designs that unify the paradigms (\cref{chap:language-and-coordination})
and the generation of macro-programs from natural language (\cref{chap:language-based-macroprogramming}).
%
Part~III asks who uses the freedom the model establishes, and how:
reconfiguration policies written by hand (\cref{chap:deployment-and-reconfiguration}),
one learned from experience (\cref{chap:heterogeneous-gnn}),
and a physical realisation of one admissible deployment (\cref{chap:demonstrator}).
%
The conclusions answer the research questions, gather the limits the chapters state, and list what is left open (\cref{chap:conclusions}).

The demonstrator closes Part~III without concluding it:
it realises one point of the deployment space on hardware and exercises none of the policies that precede it,
so the part does not culminate there.
%
A reader already familiar with the continuum and with macroprogramming can start at \cref{chap:pulverization}.
